%!TEX root=../sub_introduction.tex

Since 2004 when graphene was found by Geim and Novoselov \cite{novoselov2004electric}, the research for 2D materials is accelerated by many researchers. Graphene has layered structure in which the layers interact by weak Van der waals forces. There are many kinds of materials which have layered crystal structure like graphene. 

Among them, Transition metal dichalcogenides (TMDs), such as $MoS_2$ and $WSe_2$ are promising materials because of their special electronic properties.
Because graphene is semimetal with no band gap, its used in logic devices is strongly limited. In contrast, monolayer TMDs have a sizeable direct bandgap. Another special property is that the band structure changes dramatically with number of layers. For this reason, monolayer TMDs could be used in future nanoelectronics and optoelectronic applications.

TMDs can be enigeered by defects beyond the traditional concept of doping or alloying. This is called "Defect engineering". Starting from 2015, a range of narrow PL emissions exhibiting single-photon characteristics originating from defects on monolayers of TMDs have been detected \cite{article}. Some of these single-photon emitters are associated with excitons coupled to point defects. There are many kinds of defects in atomically thin material such as vacancy and substitution. We can intentionally make chalcogen vacancies by thermal annealing at high temperature. n this project, we hope to control the density of point defects by changing temperature of annealing. In monolayer TMDs, the exciton is very stable and the photoluminescence spectrum is mainly coming from this complex. This exction would be strongly confined in real space when coupled to defects. These localized excitons emit photon with lower energy and narrow linewidth. So, we can analyze the localized excitions by PL spectroscopy.  In summary, my objective of internship is to make defects intentionally by annealing and to analyze it by PL spectroscopy at low temperature.